\section{The Approach}\label{sec:b-approach}

The approach adds a quantum layer to srcML and srcSlice. srcML parses the program, and five extensions, E1 to E5,
give its markup a quantum meaning: they read what each call does to which qubits (E1), reconstruct the circuit the
program builds (E2), slice every qubit of that circuit (E3), add the classical statements that compute each slice (E4),
and summarize each qubit in a profile (E5). The result is shown on the program's lines. We implement the approach in a
tool, \tool{}, for Qiskit programs~\cite{javadi2024qiskit} and OpenQASM~2 circuits.

\subsection{E1, quantum roles of calls}
Qiskit applies a gate through a method of the circuit, as in \texttt{qc.cx(a, b)} or \texttt{qc.ry(theta, q)}. E1
decides, for every such call, which arguments are qubits, which are classical bits and which are angles. It reads this
from a role table derived from the signatures of Qiskit's \texttt{QuantumCircuit} methods rather than from a
hand-written list, so the table follows the library. It holds the \numToolGateQiskit{} methods of
Qiskit~\numQiskitVersion{} that take qubits and add an operation, \numToolGateLegacy{} methods that Qiskit~1.0 removed
but older programs still call, and \numToolGateSpecial{} methods whose arguments the signatures do not name uniformly,
\numToolGateTable{} in all. Table~\ref{tab:b-roles} lists them by kind, with the calls outside the table and how they are
read. A circuit composed into another is mapped onto the qubits it is given, and a gate the table does not know, such as
a circuit from Qiskit's library, becomes one operation on all its qubits, which can only enlarge a slice. A method the
approach does not model is reported to the developer, never skipped silently.

\begin{table}[t]
  \centering
  \caption{How E1 reads the calls of a Qiskit program}
  \label{tab:b-roles}
  \footnotesize\setlength{\tabcolsep}{2.6pt}\renewcommand{\arraystretch}{1.12}
  \newcommand{\sep}{\arrayrulecolor{black!25}\hline\arrayrulecolor{black}}
  \newcommand{\blk}[1]{\multicolumn{4}{@{}l}{\textit{#1}}\\}
  \begin{tabular}{@{}>{\raggedright\arraybackslash}p{2.75cm} r >{\raggedright\arraybackslash}p{2.95cm} r@{}}
    \toprule
    \textbf{Calls} & \textbf{Methods} & \textbf{Read as} & \textbf{Programs} \\
    \midrule
    \blk{In the role table}
    Single-qubit gates, \texttt{h}, \texttt{rz}, \texttt{u} & \numKindSingleN{} & one operation on its qubit & \numCallsSingle{} \\ \sep
    Two-qubit gates, \texttt{cx}, \texttt{cz}, \texttt{swap} & \numKindTwoN{} & one operation on both qubits & \numCallsTwo{} \\ \sep
    Multi-qubit gates, \texttt{ccx}, \texttt{mcx} & \numKindMultiN{} & one operation on all its qubits & \numCallsMulti{} \\ \sep
    Removed in Qiskit~1.0, \texttt{cnot}, \texttt{u3} & \numKindLegacyN{} & as gates of the same arity & \numCallsLegacy{} \\ \sep
    \texttt{measure}, \texttt{reset}, \texttt{barrier}, \ldots & \numKindSpecialN{} & a measurement writes its bit; a barrier is never sliced & \numCallsSpecial{} \\
    \midrule
    \blk{Outside the role table}
    \texttt{append}, \texttt{compose}, \texttt{measure\_all} & \numKindComposeN{} & mapped onto the target's qubits & \numCallsCompose{} \\ \sep
    \texttt{if\_test}, \texttt{while\_loop}, \texttt{switch}, \ldots & \numKindFlowN{} & conditioned, Section~\ref{sec:b-loops} & \numCallsFlow{} \\ \sep
    \texttt{inverse}, \texttt{power}, \ldots & \numKindTransformN{} & result unknown, reported & \numCallsTransform{} \\ \sep
    \texttt{circuit\_to\_dag} & \numKindDagN{} & trace of that circuit ends & \numCallsDag{} \\ \sep
    any other, \texttt{reverse\_bits}, \texttt{run}, \ldots & -- & not modelled, reported & \numUnmodMethodProgs{} \\
    \bottomrule
  \end{tabular}\\[3pt]
  \begin{minipage}{\columnwidth}\footnotesize
  Programs counts the programs of LintQ's corpus whose source calls a method of the row, and for the last row those in
  which the interpretation meets a method it does not model.
  \end{minipage}
\end{table}

\subsection{E2, the static trace}
Qiskit computes qubit indices, register sizes and angles in Python, so E2 interprets the srcML tree of the program over
the values it can know and records the operations the program would append, as a trace. Each operation carries its
source line and the iteration of every loop around it. The interpretation unrolls loops of the host program, inlines
its functions and follows its branches. Where a value cannot be known without running the program, it
over-approximates. An index it cannot evaluate acts on every qubit it may denote, both branches of a condition it cannot
decide are followed, and a loop of unknown length is repeated as Section~\ref{sec:b-loops} describes, up to
\numToolMaxCopies{} copies. Each over-approximation adds operations or qubits, so it can only enlarge a slice, and each is
reported to the developer.

\subsection{E3, the slices}
To srcSlice, \texttt{qc.cx(a, b)} defines no variable. E3 reads each operation of the trace by its qubits and bits
instead. A gate writes every qubit it acts on, its controls included, a measurement writes its bit, and a conditioned
operation reads the bits it tests. A qubit's slice is then the backward light cone of its final state, the relation
compilers use to order and place gates~\cite{murali2019triq, amy2020staq}, followed from each measurement to the
operations its bit conditions. E3 computes it in one backward pass over the trace with a set of needed qubits and bits,
in time linear in the trace, and the forward pass from an operation gives its forward slice. For large circuits, the
pass carries the set of criterion qubits on each needed location and so slices every qubit at once.

\subsection{Loops that run on the device}\label{sec:b-loops}
Since Qiskit~1.0, a circuit can contain its own loops and branches, \texttt{while\_loop}, \texttt{for\_loop},
\texttt{switch} and the \texttt{else} branch of \texttt{if\_test}, which run on the device and test measured bits. A
\texttt{for\_loop} over a known range is unrolled. Any other device loop is repeated, with every operation of its body
conditioned on the bits its condition reads, until a repetition adds no qubit or bit to the needed set.

\begin{proposition}\label{prop:b-loops}
Let a device loop with body $B$ end with probability one, let $\Lambda$ be the qubits and bits $B$ acts on, and let the
slice be computed on $|\Lambda| + 1$ conditioned copies of $B$. Deleting from $B$ every operation that none of its copies
puts in the slice of $W$ leaves the state of $W$ unchanged for every input.
\end{proposition}

\begin{IEEEproof}[Proof sketch]
A loop that runs at most $K$ times behaves as $K$ copies of $B$, each conditioned on the bits that decide the loop. In
the backward pass every copy is processed in the same way and the needed set only grows. Once a copy adds nothing, every
earlier copy keeps the same operations, and the needed set can grow at most $|\Lambda|$ times. So for every $K$ the copies
keep only operations that one of the $|\Lambda| + 1$ copies keeps. The bits that decide the loop are read by every copy,
so the measurements that write them stay in the slice and the distribution of the number of iterations is unchanged. A
loop that ends with probability one is the limit of its bounded prefixes.
\end{IEEEproof}

\subsection{E4, the classical statements}
The statements that compute a slice's qubits, angles, conditions and loop bounds come from srcSlice's profiles. E4
starts from the slice's lines and the headers of the loops, conditions and \texttt{with} blocks around them, and follows
the profiles of the variables they use back to their definitions, to the variables that flow into them, and from a
function's parameters to the arguments of its calls. srcSlice runs with a patch that adds Python and two fixes that this
work exposed, for uses inside a subscript, as in \texttt{q[a]}, which is how Qiskit names a qubit, and for module
variables read inside a function.

\subsection{E5, qubit profiles}
srcSlice records a slice profile for every variable. E5 gives each qubit a counterpart, which gathers what a developer
first asks about it.

\begin{definition}[Qubit profile]\label{def:b-profile}
The profile of a qubit $q$ records its \emph{uses}, the lines of the operations on $q$; its \emph{names}, the source
expressions that denote $q$; its \emph{functions}, the call sites of the functions in which $q$ is acted on; its
\emph{dependents}, the qubits that the first operation on $q$ reaches; its \emph{bits}, those measured from $q$, with the
lines they condition; and its \emph{backward} and \emph{forward} lines, those of its slice and of the forward slice of its
first operation.
\end{definition}

In Fig.~\ref{fig:b-example}, the profile of \texttt{q[4]} has the uses 13 to 16, 19 and 21, the names \texttt{q[b]} and
\texttt{q[src]}, the dependents \texttt{q[2]}, \texttt{q[3]} and \texttt{q[5]} to \texttt{q[8]}, and the bit \texttt{c[4]},
which conditions line 21. The two names show at once that \texttt{q[4]} is first a target and then the carrier.

\subsection{Two front ends}
Qiskit programs are read through srcML as described. OpenQASM~2 circuits, the format that Qiskit exports and benchmark
suites ship in, name their qubits explicitly and run in order, so E1 is given by the syntax and E2 reduces to reading
the statements. srcML has no released parser for OpenQASM yet, so \tool{} reads OpenQASM~2 with a reader that maps every
operation to its line and produces the same trace, and E3 to E5 apply unchanged.
